% !TEX root = ../macro.tex
\section{OLG Model}\label{sec:olg}

Ramsey assigns intertemporal choices to a dynasty that continues indefinitely.
Now separate a person's lifetime from the economy's lifetime. People work,
retire, and leave; new people arrive. Saving remains an optimal choice, but
the ownership of assets and the identity of the people paying for transfers
must be tracked by generation.

\subsection{From retirement needs to overlapping generations}

\subsubsection{An economy continues after its current savers have died}

Imagine looking at households' balance sheets at one date. Some people are
building retirement assets from wages; others are spending assets accumulated
earlier. A single consumption total conceals these opposite positions. A policy
that taxes workers and pays retirees changes the resources of different people,
even when it leaves aggregate resources unchanged at the moment of payment.

One response would be to add age labels to the Ramsey household. But if an
altruistic dynasty still pools all resources and values every descendant,
those labels need not change the underlying allocation problem. To isolate
generational turnover, instead let each person value only consumption within
his or her own finite life. Give the young labor income and the old a need to
finance consumption without working. Then ask what assets can carry resources
between these stages, and who supplies the counterpart to those claims.

Two periods are the smallest useful life cycle: one for work and saving, one
for retirement and spending. New workers arrive while earlier workers retire.
The approximation does not describe every age or predict a literal retirement
age. It makes the relation between life-cycle saving, asset ownership, and
population replacement visible with only two age groups.

\subsubsection{Choose a way to finance retirement, then close the economy}

Start with physical capital as the only asset. A young person acquires capital
claims; next period those claims pay rental income and the value of surviving
capital. Retirees consume the proceeds. A wage and a return determine an
individual saving choice, but they are not arbitrary economy-wide inputs:
workers' saving determines the next capital stock, which determines the return
those same workers receive when old. This closes the market model.

An alternative is a transfer from next period's workers. That alternative
cannot simply be assumed to be a private enforceable contract with people who
have not yet been born. A pension system or a government issuing debt adds an
institution capable of continuing across generations. We first solve the
capital-only economy, then introduce those institutions explicitly.

The production economy follows Diamond's OLG framework; the sequence of
questions above is a teaching reconstruction, not a claim about the author's
private thoughts.\footnote{Peter A. Diamond (1965), ``National Debt in a
Neoclassical Growth Model,'' \textit{American Economic Review}, 55(5),
1126--1150. Diamond's original model uses no depreciation; we keep general
$\delta$ and identify the full-depreciation course example separately.
\url{https://www.aeaweb.org/aer/top20/55.5.1126-1150.pdf}.}

\begin{blueblock}
\textbf{Meeting only once does not rule out intertemporal trade.}
A young person can save in capital and consume its payoff when old. That is
already an intertemporal allocation. The important limitation is that currently
living people cannot make ordinary private contracts with unborn generations.
Without bequests or altruistic links, no household represents the interests
of the infinite sequence of future entrants. Finite individual lives and an
infinite economic horizon can therefore have different welfare implications
from a single infinite-lived dynasty.
\end{blueblock}

% BEGIN ORIGINAL OLG COURSE ASSUMPTIONS - retained, not typeset
% \subsection{Assumptions}
% \begin{align}
% &C_{t}= L_{t} C_{1t}^{\text{young}}+L_{t-1}C_{2t}^{\text{old}} \text{ (where subscript 1 and 2 stands for age)}
% \\&L_{t}=(1+n)L_{t-1} \Rightarrow L_{t}=(1+n)^{t}L_{0}
% \\&S_{t}=s_{t}L(t)
% \end{align}
% END ORIGINAL OLG COURSE ASSUMPTIONS

\Needspace{12\baselineskip}
\subsection{Assumptions and the cohort timeline}\label{sec:olg-assumptions}

The cards state the full environment before solving it. Within a cohort people
are identical; across ages, their income, remaining lifetime, and asset position
differ. There is no single representative person for both age groups.

\begin{modelsetup}[1]
\begin{modelcondition}{01}{Dates, cohorts, and population}
Time is discrete, $t=0,1,\ldots$. A cohort of $L_t$ people is born at each date,
lives for two periods with certainty, and dies after its second period.
At date $t$, $L_t$ young coexist with $L_{t-1}$ old:
\[
L_t=NL_{t-1},\qquad N\equiv1+n>0,\qquad L_t=N^tL_0.
\]
The baseline takes $n\geq0$; the population comparison later also permits
$-1<n<0$ when the other conditions remain feasible. Each young person supplies
one unit of labor inelastically; the old supply none. Labor is $L_t$, whereas
total population is $L_t+L_{t-1}$. A model period is a life stage, not
necessarily a calendar year.
\end{modelcondition}
\begin{modelcondition}{02}{Goods, capital, and technology growth}
A closed economy produces one good, usable for consumption and capital.
There is no government, foreign trade, money, or public debt in the baseline.
Technology is fixed and normalized to one: $Y_t=F(K_t,L_t)$.
This is an explicit zero-technology-growth benchmark, so $k_t=K_t/L_t$ is
capital per worker, not per effective worker. Capital depreciates at
$0\leq\delta\leq1$ and surviving capital can be sold or consumed:
\[
Y_t+(1-\delta)K_t=C_t+K_{t+1}.
\]
General formulas retain $\delta$. The numerical and explicit full-depreciation
examples set $\delta=1$; that restriction is not needed for the general model.
\end{modelcondition}
\begin{modelcondition}{03}{Production restrictions}
Firms share a freely available technology. $F$ is continuous, twice
differentiable for positive inputs, positive there, with positive diminishing
marginal products and constant returns:
\[
F(\lambda K,\lambda L)=\lambda F(K,L),\qquad
f(k)\equiv F(k,1),\qquad Y_t/L_t=f(k_t).
\]
For $k>0$, assume
\[
\begin{gathered}
f(0)=0,\qquad f'(k)>0,\qquad f''(k)<0,\\
\lim_{k\downarrow0}f'(k)=\infty,\qquad
\lim_{k\to\infty}f'(k)=0.
\end{gathered}
\]
Cobb--Douglas, $f(k)=k^\alpha$ with $0<\alpha<1$, is a specified example,
not the only technology consistent with these conditions.
\end{modelcondition}
\begin{modelcondition}{04}{Households, ownership, and initial conditions}
There are many identical, price-taking people in each cohort. They receive
wages when young and asset income when old; firms belong to households and
any profits would accrue to their owners. No one values descendants' utility,
leaves a bequest, or receives an inheritance. The initially old are the
exception to the no-inherited-wealth statement: they enter date zero already
owning the given $K_0>0$, equally within their cohort. The young start with
labor income and no assets. Initially old consumption must be counted even
though their earlier saving decision lies outside the model's initial date.
\end{modelcondition}
\begin{modelcondition}{05}{Preferences, information, and choices}
A person born at $t$ chooses saving $a_t$ and consumption at ages one and two:
\[
U_t=u(c_{1,t})+\beta u(c_{2,t+1}),\qquad 0<\beta<1,
\]
\[
u(c)=\begin{cases}(c^{1-\theta}-1)/(1-\theta),&\theta>0,\ \theta\ne1,\\
\log c,&\theta=1.
\end{cases}
\]
There is no aggregate or individual uncertainty; the next return is correctly
anticipated. People take wages and returns as given, despite their joint
effect on equilibrium prices. There is no labor choice or liquidity shock.
Consumption is positive. The baseline interior optimum has $0<a_t<w_t$;
no borrowing or short sales of physical capital are allowed. Later pension
formulas must be checked against this asset constraint.
\end{modelcondition}
\begin{modelcondition}{06}{Firms, factor prices, and returns}
Competitive firms rent capital and hire labor at each date:
\[
\max_{K,L}\;F(K,L)-q_tK-w_tL.
\]
Interior demands imply $q_t=f'(k_t)$ and $w_t=f(k_t)-k_tf'(k_t)$.
Economic profits are zero under constant returns. The owner of a unit of
capital receives the rental payment and its remaining capital value:
\[
R_t=1-\delta+q_t,\qquad r_t=R_t-1=q_t-\delta.
\]
$R_t$ is a gross payoff, $r_t$ a net return, and $q_t$ a rental price.
Only with $\delta=1$ does $R_t=f'(k_t)$.
\end{modelcondition}
\begin{modelcondition}{07}{Budgets and age-specific consumption}
A person born at $t$ obeys
\[
c_{1,t}+a_t=w_t,\qquad c_{2,t+1}=R_{t+1}a_t.
\]
The old consume all their wealth and leave no assets after death. At date $t$,
$c_{2,t}$ is the consumption of cohort $t-1$, not the future consumption of
cohort $t$. Aggregate consumption and young people's asset purchases are
\[
C_t=L_tc_{1,t}+L_{t-1}c_{2,t},\qquad
\mathcal A_t=L_ta_t.
\]
We use $a_t$ for saving per young person to distinguish it from Solow's saving
share $s$. $\mathcal A_t$ corresponds to the original course notation for
aggregate young saving; it is not automatically national saving.
\end{modelcondition}
\begin{modelcondition}{08}{Market clearing and the equilibrium path}
Capital and labor demands equal their supplied stocks, $K_t^d=K_t$ and
$L_t^d=L_t$. Capital is the only asset, so young people's wealth at the end of
$t$ must equal next period's capital:
\[
K_{t+1}=\mathcal A_t=L_ta_t,\qquad Nk_{t+1}=a_t.
\]
Goods clearing must also hold:
\[
c_{1,t}+\frac{c_{2,t}}{N}+Nk_{t+1}
=f(k_t)+(1-\delta)k_t.
\]
An equilibrium is a path of allocations and prices satisfying household and
firm optimality, correct return expectations, and these clearing conditions,
starting from the specified initial ownership. No infinite-lived household's
transversality condition is imposed on the economy as a whole.
\end{modelcondition}
\end{modelsetup}

\subsubsection{Read a cohort's lifetime across dates}

\begin{center}
\begin{tabular}{lccc}
\hline
Cohort & Date $t-1$ & Date $t$ & Date $t+1$\\
\hline
Born at $t-1$ & Young, saves & Old, consumes & No longer alive\\
Born at $t$ & Not yet born & Young, saves & Old, consumes\\
Born at $t+1$ & Not yet born & Not yet born & Young, saves\\
\hline
\end{tabular}
\end{center}

Reading across a row gives one person's lifetime budget. Reading down a
column gives the economy's resources at one date. This distinction explains
why $c_{1,t}+c_{2,t+1}/R_{t+1}=w_t$ is an individual present-value budget,
whereas $c_{1,t}+c_{2,t}/N$ is current aggregate consumption per worker.
The two formulas combine different people and different dates.

\subsection{Household optimization and life-cycle saving}

\subsubsection{Budgets, multipliers, and the Euler condition}

Take $w_t>0$ and $R_{t+1}>0$ as given. Monotonic utility makes the two budgets
bind. With multipliers on resources remaining, write
\begin{equation}
\begin{split}
\mathscr L={}&u(c_{1,t})+\beta u(c_{2,t+1})\\
&+\lambda_{1,t}(w_t-c_{1,t}-a_t)
+\lambda_{2,t+1}(R_{t+1}a_t-c_{2,t+1}).
\end{split}
\end{equation}
At an interior optimum,
\begin{align}
u'(c_{1,t})&=\lambda_{1,t},\\
\beta u'(c_{2,t+1})&=\lambda_{2,t+1},\\
\lambda_{1,t}&=R_{t+1}\lambda_{2,t+1}.
\end{align}
Eliminating both multipliers gives
\begin{equation}\label{eq:olg-euler}
u'(c_{1,t})=\beta R_{t+1}u'(c_{2,t+1}).
\end{equation}
Giving up one consumption unit while young buys $R_{t+1}$ units when old.
The Euler condition prices that marginal transfer. The two budgets still
supply the levels. No infinite-horizon terminal condition is needed for this
person: after age two, there is neither consumption nor a bequest to choose.

\subsubsection{CRRA saving: an explicit solution and its derivatives}

Under CRRA, substitute the budgets into \eqref{eq:olg-euler}:
\begin{align}
(w_t-a_t)^{-\theta}
&=\beta R_{t+1}(R_{t+1}a_t)^{-\theta},\\
\left(\frac{w_t-a_t}{a_t}\right)^{-\theta}
&=\beta R_{t+1}^{1-\theta}.
\end{align}
Define $h(R)\equiv\beta^{-1/\theta}R^{(\theta-1)/\theta}$. Taking powers
and rearranging, step by step,
\begin{equation}
\frac{w_t-a_t}{a_t}=h(R_{t+1}),\qquad
\frac{w_t}{a_t}=1+h(R_{t+1}),
\end{equation}
so the saving function is
\begin{equation}\label{eq:olg-saving}
a_t=a(w_t,R_{t+1})
=\frac{w_t}{1+\beta^{-1/\theta}R_{t+1}^{(\theta-1)/\theta}}.
\end{equation}
The utility constant $-1/(1-\theta)$ has no effect on choices. The log case
will be obtained directly rather than substituting $\theta=1$ into that
utility expression.

Holding the return fixed, $a_w=1/(1+h)\in(0,1)$. Holding the wage fixed,
use $h'(R)=(\theta-1)h/(\theta R)$:
\begin{equation}\label{eq:olg-saving-derivatives}
a_R=-\frac{wh'(R)}{(1+h)^2}
=\frac{1-\theta}{\theta}\frac{wh}{R(1+h)^2}.
\end{equation}
Thus $a_R$ is positive for $\theta<1$, zero for $\theta=1$, and negative
for $\theta>1$. These are partial derivatives; an equilibrium price change
may also change wages and expected future returns together.

\begin{greenblock}
\textbf{A higher return need not raise retirement saving.}
A higher return makes postponing consumption more rewarding, but it also
allows a given retirement expenditure to be financed with fewer assets.
For $\theta<1$, willingness to substitute consumption across ages is strong
and saving rises. For $\theta>1$, the income effect dominates in this
retirement-only-income specification and saving falls. Under log utility,
these effects cancel in the saving share out of wages. The result depends
on the assumed income profile; it is not a universal empirical saving law.
\end{greenblock}

\subsubsection{The log case: a constant share of wages}

For $u(c)=\log c$, \eqref{eq:olg-euler} implies
\begin{equation}
\frac{c_{2,t+1}}{c_{1,t}}=\beta R_{t+1}.
\end{equation}
Substitute $c_{1,t}=w_t-a_t$ and $c_{2,t+1}=R_{t+1}a_t$:
\begin{equation}\label{eq:olg-log-saving}
R_{t+1}a_t=\beta R_{t+1}(w_t-a_t)
\quad\Longrightarrow\quad
 a_t=\kappa w_t,\qquad \kappa\equiv\frac{\beta}{1+\beta}.
\end{equation}
Saving is endogenous even though this example gives a constant share of
\emph{wages}. Solow imposes a share of total output; those are different
income bases. Log utility does not make retirement consumption independent
of the return: $c_{2,t+1}=R_{t+1}\kappa w_t$ still changes with $R_{t+1}$.

\subsection{Firm optimization and the return to ownership}

At date $t$, firms choose $K,L$ in
\begin{equation}
\max_{K,L}\;F(K,L)-q_tK-w_tL.
\end{equation}
Differentiate with respect to actual factor inputs before using constant
returns to express the prices per worker:
\begin{equation}\label{eq:olg-prices}
q_t=f'(k_t),\qquad
w_t=f(k_t)-k_tf'(k_t),\qquad
R_t=1-\delta+f'(k_t).
\end{equation}
Euler's theorem gives $Y_t=q_tK_t+w_tL_t$, so profits are zero after factor
payments. Capital owners still receive both rental income and the value of
surviving assets. That second item is an asset sale or liquidation, not
additional production in GDP.

For Cobb--Douglas,
\begin{equation}
q_t=\alpha k_t^{\alpha-1},\qquad
w_t=(1-\alpha)k_t^\alpha,\qquad
R_t=1-\delta+\alpha k_t^{\alpha-1}.
\end{equation}
A larger inherited capital intensity raises the current wage and lowers the
current return. A young saver cares about the \emph{next} return, determined
by $k_{t+1}$, because that is when the retirement payoff is received.

\begin{blueblock}
\textbf{The course formula $R=f'(k)$ is a full-depreciation case.}
With $\delta=1$, no capital survives after production, so the gross payoff
is the rental product alone. With $\delta=0$, the owner also recovers the
unit of capital and $R=1+f'(k)$. Both conventions can be useful. The
household Euler equation uses the gross payoff; replacing it by a rental
price without stating depreciation changes the saving problem.
\end{blueblock}

\subsection{Capital accumulation and equilibrium dynamics}

\subsubsection{Young wealth versus aggregate saving}

At the start of date $t$, the old own $K_t$. Since
$K_t=L_{t-1}a_{t-1}$, their aggregate consumption is
$L_{t-1}c_{2,t}=R_tK_t$; the same relation holds for the initial owners at
date zero. The young consume $L_t(w_t-a_t)$. Add them:
\begin{align}
C_t&=L_tw_t-L_ta_t+R_tK_t,\\
&=Y_t-L_ta_t+(1-\delta)K_t.
\end{align}
In the second line, substitute $R_t=q_t+1-\delta$ and
$Y_t=w_tL_t+q_tK_t$. Rearranging gives
\begin{equation}\label{eq:olg-capital-market}
K_{t+1}=L_ta_t,\qquad Nk_{t+1}=a_t.
\end{equation}
Thus the asset-clearing equation also satisfies goods clearing.

\begin{blueblock}
\textbf{Young saving is not national saving.}
The young acquire assets, while the old liquidate theirs. In this one-good
model, gross investment and net national saving are respectively
\[
I_t\equiv K_{t+1}-(1-\delta)K_t
=L_ta_t-(1-\delta)K_t,
\]
\[
S_t^{\mathrm{net}}\equiv Y_t-C_t-\delta K_t
=K_{t+1}-K_t=L_ta_t-K_t.
\]
Writing $K_{t+1}=(1-\delta)K_t+L_ta_t$ would count surviving capital twice:
young people's purchases already include acquiring that capital from the
old. With full depreciation, gross investment does equal young saving;
with general depreciation, the distinction must remain visible.
\end{blueblock}

\subsubsection{An implicit law of motion under general preferences}

Substitute the saving function and competitive prices into capital clearing:
\begin{equation}\label{eq:olg-implicit-map}
Nk_{t+1}
=a\!\left(f(k_t)-k_tf'(k_t),\ 1-\delta+f'(k_{t+1})\right).
\end{equation}
This is generally implicit in $k_{t+1}$: a larger capital stock changes the
return that induced people to accumulate it. Solving household saving at an
arbitrary fixed return is not yet a general-equilibrium solution.

At a positive stationary equilibrium,
\begin{equation}\label{eq:olg-general-steady}
Nk^*=a\!\left(f(k^*)-k^*f'(k^*),\ 1-\delta+f'(k^*)\right).
\end{equation}
Cobb--Douglas makes this condition
\begin{equation}
Nk^*=
\frac{(1-\alpha)(k^*)^\alpha}
{1+\beta^{-1/\theta}[1-\delta+\alpha(k^*)^{\alpha-1}]^{(\theta-1)/\theta}}.
\end{equation}
This scalar equation replaces the unfinished general steady-state expression
in the course notes. A numerical root may be needed; the explicit log result
below is a separate tractable case.

To see why stability needs conditions, differentiate
\eqref{eq:olg-implicit-map}. Since $w'(k)=-kf''(k)>0$,
\begin{equation}\label{eq:olg-map-slope}
\left[N-a_R f''(k_{t+1})\right]\frac{dk_{t+1}}{dk_t}
=a_w[-k_tf''(k_t)].
\end{equation}
If the bracket is nonzero, this gives the local slope. A locally defined map
is stable at its steady state when the absolute slope is below one.
For $\theta\leq1$, $a_R\geq0$ makes the bracket positive. For $\theta>1$,
the return response can overturn that sign or make the equation locally
singular. The log example's uniqueness and convergence cannot be asserted
for every OLG specification.

\subsubsection{Log utility and Cobb--Douglas: solve the whole path}

Using \eqref{eq:olg-log-saving}, capital clearing reduces to
\begin{equation}\label{eq:olg-log-map}
k_{t+1}=\frac{\kappa(1-\alpha)}{N}k_t^\alpha
\equiv M k_t^\alpha,
\qquad M\equiv\frac{\kappa(1-\alpha)}{N}>0.
\end{equation}
This formula holds for general $\delta$: log saving does not respond to the
return, and current wages do not depend on depreciation. Depreciation still
changes the old's consumption and welfare.

At a positive steady state, divide $k^*=M(k^*)^\alpha$ by $(k^*)^\alpha$:
\begin{equation}\label{eq:olg-log-steady}
(k^*)^{1-\alpha}=M,\qquad
k^*=\left[\frac{\beta(1-\alpha)}{(1+\beta)N}\right]^{1/(1-\alpha)}.
\end{equation}
The map is increasing; $k_{t+1}/k_t=Mk_t^{\alpha-1}$ exceeds one below $k^*$
and is below one above it. It maps either side toward the steady state without
crossing it. Thus every $k_0>0$ converges monotonically; zero is an absorbing
boundary, not the initial condition of the positive-capital result.

There is also an exact solution in log deviations. Divide the map by its
steady-state version and take logarithms:
\begin{equation}
\log(k_{t+1}/k^*)=\alpha\log(k_t/k^*),\qquad
k_t=k^*\left(\frac{k_0}{k^*}\right)^{\alpha^t}.
\end{equation}
The local slope is $\alpha\in(0,1)$. Capital-level deviations satisfy
$k_t-k^*\approx\alpha^t(k_0-k^*)$ only locally; log deviations decay exactly
at that factor. This separates convergence of a path from efficiency of the
allocation to which it converges.

\begin{greenblock}
\textbf{What the log example does and does not establish.}
Young people optimally save a fixed share of wages. Wages depend on the
inherited stock; their asset purchases become the next stock. This gives a
unique positive steady state and a stable path in this example. It does not
show that all OLG economies converge, that capital is socially optimal, or
that a household's constant saving share equals Solow's imposed output share.
\end{greenblock}

\subsection{Comparing Solow, Ramsey, and OLG}

Use a common discrete-time benchmark here: fixed technology, population
factor $N$, and depreciation $\delta$. Let $\sigma$ denote Solow's imposed
saving share to keep it distinct from individual saving $a_t$.

\begin{center}
\begin{tabular}{p{0.14\linewidth}p{0.24\linewidth}p{0.25\linewidth}p{0.25\linewidth}}
\hline
 & Solow & Ramsey dynasty & Two-period OLG\\
\hline
Saving & Given share of output & Chosen for a continuing dynasty & Chosen for each cohort's two ages\\
People & No age structure in the saving rule & Descendants linked by the objective & New cohorts enter; no altruistic link\\
Main closure & Accumulation equation & Accumulation, Euler, and terminal conditions & Cohort budgets and asset clearing\\
Efficiency & Not guaranteed by the saving rule & First-best under its stated assumptions & Overaccumulation is possible\\
\hline
\end{tabular}
\end{center}

The corresponding aggregate resource and accumulation conditions are
\begin{align}
\text{Solow:}\quad
Nk_{t+1}&=(1-\delta)k_t+\sigma f(k_t),\\
\text{Ramsey:}\quad
Nk_{t+1}&=f(k_t)+(1-\delta)k_t-c_t,\\
\text{OLG:}\quad
Nk_{t+1}&=a(w_t,R_{t+1})\\
&=f(k_t)+(1-\delta)k_t-c_{1,t}-c_{2,t}/N.
\end{align}
In the Ramsey row, $c_t$ is aggregate dynasty consumption per worker. Its
population-weighted objective in this discrete convention gives the
dynasty Euler equation $u'(c_t)=\beta R_{t+1}u'(c_{t+1})$, together with an
appropriate infinite-horizon terminal condition. The OLG Euler equation
instead links two ages of one person. The equations look similar but have
different coverage of future generations.

\begin{blueblock}
\textbf{The common resource constraint does not make the models identical.}
In OLG the old spend asset income and the young buy assets, so clearing gives
$Nk_{t+1}=a_t$. In Ramsey the continuing household chooses its next wealth
while valuing its members' future consumption. Solow specifies saving without
that optimization. The distinction is who makes the choice and whose future
is included, not whether one model respects accounting and another does not.
\end{blueblock}

\subsection{Dynamic inefficiency}
Because of the micro-foundations (patience denoted as $\beta=\frac{1}{1+\rho}$) in Ramsey, there will be no over-accumulation of capital, avoiding dynamic inefficiency. Whereas in Solow model the inefficiency is inevitable. Somehow in OLG model, the dynamic inefficiency returns.

\textbf{Solow}

The model has key equations
\begin{equation}
\begin{cases}
  Y=F(K,L)
  \\
  I=S=sY
  \\
  L_{t+1}=(1+n)L_t
  \\
  k_{t+1}=\frac{(1-\delta)k_t+sf(k_t)}{(1+n)}\iff \dot k=sf(k)-(n+\delta)k
\end{cases}
\end{equation}

For a certain production $y=k^a$ at steady state

use $k_{t+1}=\frac{(1-\delta)k_{t}+s f(k_{t})}{1+n}$
\begin{align}
&\Rightarrow(1+n)k^*=(1-\delta)k^*+s f(k^*)
\\&-(n+\delta)k^*=-s f(k^*)
\\&f(k^*)-(n+\delta)k^*=(1-s) f(k^*)=c^*
\\&\Rightarrow c^*=f(k^*)-(n+\delta)k^*
\\&\text{To maximize welfare at steady state} \iff \max c* \iff FOC: k^*_g=(\frac{n+\delta}{\alpha})^{1-\alpha}
\end{align}
also at steady state
\begin{align}
&(1+n)k^*=(1-\delta)k^*+s f(k^*)
\\&\Rightarrow k^*=(\frac{n+\delta}{s})^{1-\alpha}
\end{align}

for certain $\alpha$ and $s$ it could be
\begin{align}
\Rightarrow k^*<k^*_{g}
\end{align}
meaning there is possible dynamic inefficiency in Solow model

\textbf{Ramsey}

Ramsey has no dynamic inefficiency, here's why

The model has three key equations
\begin{align}
&\begin{cases}L_{t}=(1+n)L_{t-1} \\
k_{t+1}-(1-\delta)k_t=f(k_t)-c_t \\
\frac{u'(c_{t})}{u'(c_{t+1})}=\beta[f'(k_{t+1})+(1-\delta)]
\end{cases}
\end{align}

Put $k_{t+1}-(1-\delta)k_t=f(k_t)-c_t$ at steady state
\begin{align}
&c^*=f(k^*)-k^*+(1-\delta)k^*=f(k^*)-\delta k^*
\\& \text{To maxmize welfare at steady state}\Rightarrow \text{FOC: } f'(k^{*}_{G})=\delta
\end{align}

Put $\frac{u'(c_{t})}{u'(c_{t+1})}=\beta[f'(k_{t+1})+(1-\delta)]$ at steady state resulting
\begin{align}
&\Rightarrow 1= \beta [f'(k^{*})+(1-\delta)] =\frac{1}{1+\rho} [f'(k^{*})+(1-\delta)]
\\&\Rightarrow 1+\rho=f'(k^{*})+1- \delta
\\&\Rightarrow f'(k^{*})=\rho +\delta < f'(k^*_g)
\\&\Rightarrow k^{*}_{G}>k^{*}
\end{align}
Capital per capita at steady state is lower than $k_{gold}$, meaning it is dynamic efficient.

\textbf{OLG}

OLG in central planner form:
\begin{align}
\begin{split}
&\max \sum\limits_{t=0}^{\infty}(\beta_{s})^{t}L_{t}[u(c_{1t})+\beta u(c_{2,t+1})] \text{ ($\beta_s$ is the discount factor)}
\\&s.t. \;
c_{1t}+\frac{c_{2t}}{1+n}+(1+n)k_{t+1}=f(k_{t})
\end{split}
\end{align}

Equation of motion of capital per capita at the steady states
\begin{align}
&c^{*}\equiv c_{1}^{*}+\frac{c_{2}^{*}}{1+n}=f(k^{*})-(1+n)k^{*}
\end{align}

Since we are maximizing welfare at steady state, let $\frac{\partial  c^{*}}{\partial  k^{*}}=0$ to acquire first order condition
\begin{align}
&\Rightarrow f'(k^{*}_{G})=1+n
\end{align}
using production function $y=k^\alpha$ at equilibrium
\begin{align}
f'(k^{*})&=\alpha (k^{*})^{\alpha-1}
\\&=\alpha\{[\frac{\beta(1-\alpha)}{(1+\beta)(1+n)}]^{\frac{1}{1-\alpha}}\}^{\alpha-1}
\\&=\frac{\alpha}{1-\alpha}\frac{1+\beta}{\beta}(1+n)
\end{align}
for certain value of $\alpha$
\begin{align}
&f'(k^{*})<1+n
\\&\Rightarrow f'(k^{*})<f'(k^{*}_{G})
\\& \Rightarrow k^{*}>k^{*}_{G}
\end{align}
Meaning there can be over-accumulation of capital $\iff$ suffices dynamic inefficiency possible in OLG model.

\textbf{Conclusion}
\begin{itemize}
\tightlist
  \item Due to ad hoc given $s$, Solow model is unable to assure dynamic efficiency.
  \item Due to endogenous saving decision and discount factor $\beta$, Ramsey model is able to achieve dynamic efficiency.
  \item Due to limited life span of individuals, OLG model fails to achieve dynamic efficiency.
\end{itemize}

\subsection[Population, pensions, and public debt (planned)]{Population, pensions, and public debt}
\textit{Planned section.} Extend the two-generation benchmark to population change, pension arrangements,
and government debt. State the conditions for dynamic inefficiency before
comparing intergenerational transfers.
